% !TITLE 生年不满百
% !AUTHOR 无名氏
% !DYNASTY 两汉
% !GENRE 五言古诗
% !ORDER 2

\begin{poem}
生年不满百，常怀千岁忧\textsuperscript{1}。
昼短苦夜长，何不秉烛游\textsuperscript{2}？
为乐当及时，何能待来兹\textsuperscript{3}？
愚者爱惜费\textsuperscript{4}，但为后世嗤。
仙人王子乔\textsuperscript{5}，难可与等期\textsuperscript{6}。
\end{poem}

\begin{annotations}
\notetext{1}{千岁忧：千年的忧虑。形容人们为久远的未来、无穷的事物而徒增的世俗烦恼。}
\notetext{2}{秉烛游：拿着蜡烛在夜里游玩。秉，拿着。形容珍惜光阴、及时寻找快乐。}
\notetext{3}{来兹：来年，明年。兹，此，指代年岁。}
\notetext{4}{爱惜费：吝啬钱财，舍不得花费。费，钱财。}
\notetext{5}{王子乔：传说中的仙人。相传他是周灵王的太子姬晋，喜欢吹笙，后随道士去往嵩山，乘白鹤成仙而去。}
\notetext{6}{等期：期待达到同等的寿命。等，同等；期，期望的寿命。指普通人无法像仙人那样追求长生不死。}
\end{annotations}

\begin{appreciation}
这首诗出自古诗十九首，是一首充满生命焦虑的“及时行乐”之歌。

“生年不满百，常怀千岁忧”——人活不到一百岁，却常常怀着千年的忧虑。这是何等的矛盾！我们的生命如此短暂，却把大量的时间用来担心未来、计较得失。这种“千岁忧”，可能是功名利禄，可能是子孙后事，也可能是对死亡的恐惧。但无论如何，用有限的生命去承担无限的焦虑，本身就是一种荒谬。

“昼短苦夜长，何不秉烛游？”——白天太短，夜晚太长，为什么不拿着蜡烛继续游乐呢？“昼短苦夜长”的“苦”字很妙，既是苦于白天太短，也是苦于夜晚太长。时间对每个人都是公平的，但感受却因人而异：快乐的人觉得白天太短，忧愁的人觉得夜晚太长。诗人的建议是：既然白天不够用，那就点燃蜡烛，把夜晚也变成白天。

“为乐当及时，何能待来兹”——快乐要及时，怎么能等到明年呢？“来兹”就是来年。这种“及时行乐”的思想，在东汉末年社会动荡的背景下，有它特殊的时代意义。当时政治腐败、战乱频仍，人们看不到未来的希望，于是把目光转向当下，在有限的时光里寻求生命的快感。

“愚者爱惜费，但为后世嗤”——吝啬的人爱惜钱财，只会被后人嘲笑。“仙人王子乔，难可与等期”——成仙的王子乔，也不是我等可以企及的。这两句是对前面“及时行乐”观点的论证：既不能吝啬地活着，又不能指望长生不死，那么唯一合理的选择，就是珍惜眼前的时光。

这首诗表面上是享乐主义的宣言，骨子里却是对生命短促的深沉悲叹。它不是教人堕落，而是在有限与无限的矛盾中，试图找到一种自洽的生活方式。这种对生命意义的追问，贯穿了中国文学史，从《古诗十九首》到曹操的“对酒当歌，人生几何”，再到李白的“人生得意须尽欢”，都是同一个主题的不同变奏。
\end{appreciation}
